\section{Parser-Combitators for Path Querying}

Parser-Combinators is a way to describe context-free grammar in terms of functions and operations on them. Parser is a function which takes some input and returns either Success with result of parsing or Error in a case of failure. Such parsers are composable which makes it easy to create new parsers from existing ones.

One possible solutions of solution to create queries using parser-combinators is a Trails~\cite{ScalaGraphParsing}. But Trails strugles with left-recursion grammars like presented in Fig.~\ref{fig:exampleGrammar} and also may stuck on some graphs with loops by not yielding some paths in result stream.

Our work based on Meerkat parser-combinators library which can hadnle left recursion and as result has Shared Packed Parse Forest SPPF~\cite{SPPF} which is a graph stores all possible ways to parse given input. From that SPPF we can get everything we need to know about our paths.

But Meerkat was made to work on a linear input. And we extend input for Meerkat from linear to the graph one. That allows us to get all possible paths in graphs which is described by grammar.

Let us introduce an example. Let's assume we have context-free grammar $G$ presented on Fig. ~\ref{fig:exampleGrammar}. It produces a language whcih consists of of words starting with $b$ or $c$ and followed by some $a$'s. Its Meerkat representation is in Fig. ~\ref{fig:exampleGrammarMeerkat}

\begin{figure}[h]
\begin{align*}
& A \rightarrow\ A\ a\ |\ B \\
& B \rightarrow b
\end{align*}
\caption{Example grammar}
\label{fig:exampleGrammar}
\end{figure}


\begin{figure}[h]
\begin{lstlisting}
val A = syn(A ~~ "a" | B)
val B = syn("b" | "c") 
\end{lstlisting}
\caption{Example grammar in terms of parser combinators}
\label{fig:exampleGrammarMeerkat}
\end{figure}

Let's closely take a look at it. For every nonterminal in our CF grammar we create a val of  \lstinline{Nonterminal} type. \lstinline{syn} is a macro which automaticaly assigns a name of our val to our nonterminal. Inside a \lstinline{syn} macro we've got a defenition of nonterminal. It using two combinators \lstinline{~~} and \lstinline{|}. The first one \lstinline{~~} says that after edge described by left operand we would like to have adjacence edge described by right one. \lstinline{|} is an altenative combinator which have lower priority then \lstinline{~~} and used to describes possible alternativesof paths description. A new combinators can be created using existing ones like for nonterminal on Fig.~\ref{fig:exampleGrammarMeerkat} which makes parser combinators a powerfull technique for describing a paths in graphs.


Let's parse graph from Fig.~\ref{fig:graph}

\begin{figure}[h]
\includegraphics{graph}
\caption{Example graph}
\label{fig:graph}
\end{figure}

To do it let us use meerkat function \lstinline{parseGraphFromAllPositions(parser,graph)} which applies given parser to given graph and gets from SPPF all pairs of nodes such that exists a path between them described by a parser.

The result for graph in Fig.~\ref{fig:graph} is $\{(3,2), (0,2), (3,1), (0,1)\}$, where $(i,j)$ stands for the path from node with label $i$ to the node with label $j$


